% bank_fintech_tone/drafts/en/manuscript.tex
% P4.2 manuscript skeleton (plan-r3 derivative-paper extension, 2026-05-15)
% Compile: xelatex + biber + xelatex × 2

% papers_assets/_shared/preamble_en.tex
% Shared English LaTeX preamble for all 5 paper drafts.
% P4.1 (plan-r3 derivative-paper extension, 2026-05-15)
% Compile: xelatex + biber + xelatex × 2

\documentclass[12pt,a4paper]{article}

% --- Fonts & Encoding ---
\usepackage{fontspec}
\setmainfont{Times New Roman}[Ligatures=TeX]

% --- Math ---
\usepackage{amsmath, amssymb, amsthm}

% --- Tables & Figures ---
\usepackage{booktabs}
\usepackage{graphicx}
\usepackage{longtable}
\usepackage{rotating}
\usepackage{caption}
\captionsetup{font=small,labelfont=bf,skip=4pt}

% --- Layout ---
\usepackage[margin=1in]{geometry}
\usepackage{setspace}
\onehalfspacing
\usepackage{hyperref}
\hypersetup{colorlinks=true,linkcolor=blue,citecolor=blue,urlcolor=blue}

% --- Bibliography (biblatex + biber) ---
\usepackage[backend=biber,style=apa,natbib=true]{biblatex}

% --- Useful macros ---
\newcommand{\code}[1]{\texttt{#1}}
\newcommand{\NoteThis}[1]{\textbf{\textcolor{red}{[TODO: #1]}}}
\usepackage[dvipsnames]{xcolor}

% --- Theorem environments ---
\newtheorem{prop}{Proposition}
\newtheorem{hyp}{Hypothesis}


\addbibresource{refs.bib}

\title{Media Tone, FinTech Strategy, and Bank Performance:\\
       Evidence from Chinese Bank--Year News Coverage}
\author{\NoteThis{Author names}}
\date{\today}

\begin{document}
\maketitle

\begin{abstract}
\NoteThis{Final abstract pending.}
We construct a bank--year panel from approximately 37~GB of Chinese
media reports on financial technology (FinTech) and study the
relationship between media tone toward bank FinTech strategy and bank
outcomes. Using a three-family large language model (LLM) ensemble for
sentiment annotation, we measure both the locus (full-text vs.\
entity-context) and the substantiveness (full-text vs.\ entity-context)
of FinTech coverage. We find \NoteThis{key result placeholder}\@.
\end{abstract}

\section{Introduction}\label{sec:intro}
\NoteThis{Motivate the puzzle: media coverage of bank FinTech strategy
matters for institutional theory (Westphal \& Zajac 1994) and for digital
transformation outcomes (Gomber et al.\ 2018, \emph{JMIS}). Reference
Hassan et al.\ (2019, \emph{QJE}) firm-level political-risk methodology
as our text-measurement template.}

\section{Literature Review}\label{sec:lit}
\NoteThis{Three strands: (1)~media-as-information-intermediary
(Tetlock 2007); (2)~symbolic vs.\ substantive corporate response
(Westphal--Zajac 1994/1998 \emph{ASQ}); (3)~Bartik instrument design
(Goldsmith-Pinkham et al.\ 2020, \emph{AER}).}

\section{Methodology}\label{sec:method}
\subsection{Sample and Data}
The bank--year panel covers \NoteThis{N=164 bank--years over the
2017--2018 period (extend after additional vintages merged)}. Raw text
inputs are XLSX media reports processed via the \code{news\_process}
pipeline (Phase~9b baseline).

\subsection{Tone Annotation (F1: Three-Family LLM Validation)}
We employ three independent LLM families for sentiment annotation:
\code{gpt-5.4-nano} (OpenAI), \code{gemini-3.1-flash-lite}
(Google), and \code{deepseek-v4-flash} (DeepSeek). The OpenAI and
Google models are routed through the yunwu.ai aggregator using the
OpenAI-compatible \code{/v1/chat/completions} API; the DeepSeek model
is called directly at \code{api.deepseek.com}. All calls use
\code{temperature=0} and \code{seed=42} with a byte-frozen prompt
(version~v8) so that outputs are deterministic and cache-reproducible.
Inter-family agreement is reported as Cohen's $\kappa$ (labels
remapped to non-negative codes before computation) and Krippendorff's
$\alpha$ (nominal, missing-label-aware, tolerant of two- vs.\
three-class label sets); majority vote across families produces the
final sentiment label, with ties marked as \code{undefined} (not
\code{neutral}).

Beyond the majority-vote consensus (\code{llm\_vote\_tone\_full}, the
headline treatment per the paper specification), we retain each
family's own tone series
(\code{gpt54nano\_tone}, \code{gemini31\_tone},
\code{deepseekv4\_tone}) as \emph{co-equal mutual-robustness}
measures: the three per-model series are reported side by side, each
in its own just-identified specification, and are never pooled into a
single over-identified system (instruments derived from the same
articles share common shocks, rendering a joint Hansen $J$
uninformative). Because the per-model series differ in raw variance,
cross-model coefficient comparisons standardize each tone measure to
$z$-scores. Family-wise error across the per-model specifications is
controlled by Romano--Wolf step-down resampling
(\code{rwolf}; Table~A12), rather than by informal
``all-three-significant'' narration.

\subsection{Article-Level Coding (F2: Decontamination)}
Articles where (a)~all three LLM families fail to return a valid response
(\code{llm\_call\_failure=1}) or (b)~the three sentiment votes split
three ways with no majority (\code{llm\_undefined=1}) are coded as
\emph{missing} rather than neutral. Articles where the bank--FinTech
relevance fails majority verification on the full-text track
(\code{llm\_verified\_full=False})
likewise have their sentiment treated as missing. This is a deliberate
departure from prior versions of the panel which coerced these to
``neutral'' and conflated 81.7\,\% of the neutral labels with measurement
failure or verification rejection.

\paragraph{Full-text-only retention gate.}
Whether an article counts as bank--FinTech coverage at all is decided
solely on the full-text track (\code{llm\_verified\_full}). The
entity-context (\code{ctx}) and same-sentence co-occurrence
(\code{ss}) tracks are \emph{conditional-on-full} measures: their row
universe is a subset of the full-track retained set, and they serve as
measurement-window robustness tracks rather than independent retention
gates. Relative to earlier panel vintages that retained an article if
\emph{any} track verified it, this stricter gate mechanically shrinks
article counts and derived shares (\code{n\_articles},
\code{fintech\_share}, \code{attention\_pct}). These level shifts are
definitional consequences of aligning the retention gate with the
headline treatment track---not data errors---and the direction of the
change is conservative (it removes marginal articles that only a
narrower window verified).

\subsection{Zero-Observation Semantics (F3: NaN Main Path)}
For bank--years with zero FinTech news coverage (\code{n\_articles=0}),
cosine-similarity-based features (\code{auth\_sim\_*},
\code{benefit\_*}, \code{article\_sim\_*}), ratios
(\code{token\_adopt\_*}, \code{hidden\_adopt\_*},
\code{subst\_adopt\_*}), and continuous intensity scores
(\code{fintech\_intensity\_mean}) are coded as undefined (\code{NaN})
rather than zero. We follow Allison (2001) and Heckman (1979) in treating
these as missing-by-design and use listwise deletion in regressions.
Count and sum-type variables remain at~$0$ as the mathematically correct
zero-observation value.

\subsection{Bartik IV Pruning (F4)}
Of approximately 30 candidate Bartik instruments in earlier panel
vintages, 28 are removed in v6 for violating Goldsmith-Pinkham et al.\
(2020, AER) exogeneity. The retained instruments are
\code{bartik\_type\_iv\_full} and \code{prov\_leave1out\_tone\_full}.

\subsection{Embedding-Based Semantic Measures (Co-Primary Dual Track; F5)}
All embedding-derived features (FinTech intensity, policy alignment,
semantic staleness, and related similarity measures) are computed
under \emph{two co-primary embedding models from independent vendors},
reported side by side as mutual robustness: Alibaba
\code{Qwen3-Embedding-8B} (4{,}096 dimensions, accessed via the
SiliconFlow API; canonical, unsuffixed columns) and OpenAI
\code{text-embedding-3-large} (3{,}072 dimensions; \code{\_lg}-suffixed
columns). The two tracks obey strict column parity (every embedding
feature exists in both tracks or in neither), each track enters its
own just-identified specification, and the two are never combined into
a single over-identified system. Cross-vendor independence
(Alibaba vs.\ OpenAI) guards against vendor-specific representation
artifacts driving any single result.

\emph{Vendor-change disclosure.} Earlier pipeline vintages used a
Google Gemini embedding model in the canonical slot. As of pipeline
v3.15.0 (June 2026) the canonical slot is Alibaba
\code{Qwen3-Embedding-8B}; all embedding-derived features were
regenerated under a controlled full re-embedding, and embedding-space
semantic validity is gated in the test suite by a cross-vendor
Spearman rank-correlation check between the two co-primary embedding
spaces ($\rho \geq 0.80$, $\geq 500$ pairs). The sentiment-annotation LLM ensemble
(Section~\ref{sec:method}) is unaffected by this change: embedding
models and sentiment LLMs are disjoint components of the pipeline.

\section{Data}\label{sec:data}
\NoteThis{Descriptive statistics in Table~1
(\code{news\_process/papers/bank\_fintech\_tone/tables.py:build\_table1}).}

\section{Main Results}\label{sec:main}
\NoteThis{Table 2 (M1--M6 progressive controls) and Figure 2 (coefplot)
from \code{output/tables/table2\_main\_m1m6.csv}.}

\section{Mechanism}\label{sec:mech}
\NoteThis{Table 3 (mechanism: technology-application vs.\ substantive--
symbolic gap) from \code{output/tables/table3\_mechanism.csv}.}

\section{Heterogeneity}\label{sec:het}
\NoteThis{Table 4 (bank size $\times$ frontier $\times$ geography) from
\code{output/tables/table4\_heterogeneity.csv}.}

\section{Robustness}\label{sec:robust}
\NoteThis{Table 5 (placebo, leave-one-out, wild bootstrap) plus the
1016-line falsification template
(\code{papers\_assets/bank\_fintech\_tone/do/falsification\_template.do.tmpl}).}

Cross-model robustness: the headline specification is re-estimated
separately with each LLM family's own tone measure (and, analogously,
with each of the three Chinese financial-domain BERT models'
signed net-tone measures), with per-model coefficients standardized to
$z$-scores and family-wise error controlled by Romano--Wolf step-down
resampling (Table~A12). Inter-model agreement diagnostics
(Cohen's $\kappa$ on remapped non-negative codes; Krippendorff's
$\alpha$, missing-label-aware) accompany the table.

\section{Conclusion}\label{sec:concl}
\NoteThis{Final discussion and policy implications.}

\printbibliography

\end{document}
